(1) Under (S) the states of Proposition \claimref{standard-model.quarks-leptons-two-level-fermion-content}(1) and (2) have total mode numbers $2$, $1$ and $4$ for $(1,1)$, $(1,0)$ and $(1,3)$, and by (S) the chirality is the parity; so $(1,1)$ and $(1,3)$ share a handedness and $(1,0)$ has the other. By Fact \claimref{standard-model.imported-anomaly-coefficients}'s convention a state of the opposite handedness is listed as its conjugate, so the listing that keeps the component $(1,1)$ keeps the component $(1,3)$ and conjugates the component $(1,0)$ to $(2,1)_{-1/2}$, the same representation as the component $(1,3)$ (Proposition \claimref{standard-model.chirality-two-halves-spinor-module-conjugation}). The $[SU(2)]^{2}U(1)$ coefficient is the sum of $Y$ over the doublets weighted by colour multiplicity: $3 \cdot \tfrac16 + (-\tfrac12) = 0$ for the listing above, and $3 \cdot \tfrac16 + \tfrac12 = 1$ for the listing that keeps $(1,0)$ and conjugates $(1,3)$.

(2) Proposition \claimref{standard-model.quarks-leptons-two-level-fermion-content}(3) gives the single-twistor states as $(1)_{0}$ at grade $0$, $(3)_{-1/3}$ at grade $1$, $(\bar{3})_{-2/3}$ at grade $2$ and $(1)_{-1}$ at grade $3$. The even grades lie in $\T$ and the odd grades in $\T'$, the two chirality halves of Fact \claimref{standard-model.imported-pair-systems-two-chirality-halves}. Conjugating the odd ones, with the conjugate representation and the opposite hypercharge: $(3)_{-1/3} \mapsto (\bar{3})_{+1/3}$ and $(1)_{-1} \mapsto (1)_{+1}$. The even grades are unchanged. This is the display. Under (S) the grade-$k$ single-twistor state is the $(0, k)$ component, of total mode number $k$, so the parity used here is the parity of (1).

(3) Consistency. $[SU(3)]^{2}U(1)$: the triplets and antitriplets are $Q$ (multiplicity two) and $u^{c}$, $d^{c}$: $2 \cdot \tfrac16 - \tfrac23 + \tfrac13 = 0$. $[SU(2)]^{2}U(1)$: (1). Gravitational: $6 \cdot \tfrac16 + 2 \cdot (-\tfrac12) + 3 \cdot (-\tfrac23) + 3 \cdot \tfrac13 + 1 + 0 = 1 - 1 - 2 + 1 + 1 = 0$. Cubic: $6 \cdot \tfrac{1}{216} + 2 \cdot (-\tfrac18) + 3 \cdot (-\tfrac{8}{27}) + 3 \cdot \tfrac{1}{27} + 1 + 0 = \tfrac{1}{36} - \tfrac14 - \tfrac89 + \tfrac19 + 1 = 0$. Uniqueness. Write $Y_{Q} = a = \tfrac16$, $Y(u^{c}) = b$, $Y(d^{c}) = c$, $Y_{L} = d$, $Y(e^{c}) = e$, and $f$ for the fourth singlet, with $f = 0$ or absent. $[SU(2)]^{2}U(1)$: $3a + d = 0$, so $d = -\tfrac12$. $[SU(3)]^{2}U(1)$: $2a + b + c = 0$, so $b + c = -\tfrac13$. Gravitational: $6a + 3b + 3c + 2d + e + f = 1 - 1 - 1 + e = 0$, so $e = 1$. Cubic: $6a^{3} + 3(b^{3} + c^{3}) + 2d^{3} + e^{3} + f^{3} = \tfrac{1}{36} - \tfrac14 + 3(b^{3} + c^{3}) + 1 = 0$, so $b^{3} + c^{3} = -\tfrac{7}{27}$; with $s = b + c = -\tfrac13$ and $p = bc$, the identity $b^{3} + c^{3} = s^{3} - 3ps$ gives $-\tfrac{1}{27} + p = -\tfrac{7}{27}$, so $p = -\tfrac29$, and $b, c$ are the roots of $y^{2} + \tfrac13 y - \tfrac29 = 0$, namely $(-\tfrac13 \pm 1)/2 = \{+\tfrac13, -\tfrac23\}$. Unique. Necessity of the fourth-singlet hypothesis: if $f$ is free, the gravitational and cubic conditions constrain $e + f$ and $e^{3} + f^{3}$ together and no longer fix $e$. Explicitly, $(a, b, c, d, e, f) = (\tfrac16, -\tfrac16, -\tfrac16, -\tfrac12, \tfrac12, \tfrac12)$ satisfies all four conditions ($b + c = -\tfrac13$; $d = -\tfrac12$; $1 - 1 - 1 + \tfrac12 + \tfrac12 = 0$; $\tfrac{1}{36} - \tfrac{1}{36} - \tfrac14 + \tfrac18 + \tfrac18 = 0$), and it is not isolated: $Y + t(B - L)$, with $B - L$ the baryon-minus-lepton charge ($\tfrac13$ on $Q$, $-\tfrac13$ on $u^{c}$, $d^{c}$, $-1$ on $L$, $1$ on $e^{c}$ and on the fourth singlet), solves the four conditions for every $t$: each of the four sums is a polynomial of degree at most three in $t$, so vanishing at four values of $t$ is vanishing identically; the script confirms the vanishing at $t = 0$ (the derived spectrum), at $t = 1, -1, 2$, and at $t = \tfrac57$. (Rescaling a member to $Y_{Q} = \tfrac16$, possible for $t \neq -\tfrac12$ where $Y_{Q} = \tfrac16 + \tfrac{t}{3}$ vanishes, multiplies the three linear sums by a constant and the cubic one by its cube, so preserves vanishing; the argument is run on the unrescaled family.) The exhibited point is the $t \to \infty$ limit of the family, rescaled. With $f = 0$ imposed, or the fourth singlet absent, the family collapses to the single solution above. Necessity of the common handedness: conjugating the two doublets alone sends $(a, d)$ to $(-\tfrac16, \tfrac12)$ in the conjugate representations, and the four sums become $[SU(3)]^{2}U(1) = -2a + b + c = -\tfrac23$, $[SU(2)]^{2}U(1) = -3a - d = 0$, gravitational $-6a + 3b + 3c - 2d + e = 0$, cubic $-6a^{3} + 3(b^{3} + c^{3}) - 2d^{3} + e^{3} = +\tfrac49$. Every line of this arithmetic is verified exactly by the script.
